\section{The two constant terms of a three-letter mixed arc}

% BEGIN K83J1
\paragraph{K83J1. An integral constant at the outer primitive level.}
Take the actual mixed arc of U83M1, read forward from $c-2$ to
$c$, with signs down, up, down and length three. Write
\[
 y=s_{c-2}+S,\quad b_1=s_{c-1},\quad B=s_c+S,\quad
 A=O[[v]],\quad P_j=\varphi^j(v+p).
\]
Here $c-2$ is nonthick and the other two positions are thick.
Suppose $b_1<p(y+\eta)$, the near case of the native
length-three dichotomy. Fix any scalar $f$ of S83M1.
In the indicated component its conditions say that $P_1^{b_1}f\in A$
and that $f$ is congruent to constants $x_2,x_0\in L$ at the
$P_2$ and $P_0$ levels, respectively, to precisions $y,B$ after
multiplying by $P_1^{b_1}$ and inverting $p$. Define
\[
 K=0\quad(b_1<py),\qquad
 K=d-y+1+v_p\binom d{y-1},\quad
 d=\lceil b_1/p\rceil\quad(b_1\ge py).
\]
Then $X_2=p^K x_2\in O$ and
\[
 p^Kf=X_2+f_2,\qquad P_1^{b_1}f_2\in P_2^y A .
\]
The residual still has a native constant at $P_0$: for some
$\kappa\ge0$, $X_0\in O$,
$P_1^{b_1}(p^\kappa f_2-X_0)\in P_0^B A$.

To see the integral cost, in the finite free $O$-module
$A/(P_2^y)$ the scalar multiple $x_2P_1^{b_1}$ equals the
image of the integral numerator of $f$. Its content is the
minimum $p$-valuation of coordinates in any integral basis.
We claim the content of $P_1^{b_1}$ is at most $K$.
If $b_1<py$, reduction modulo the coefficient uniformizer gives
the nonzero class of a unit times $v^{pb_1}$ modulo $v^{p^2y}$,
so the content is zero.

For the other case put $Y=P_1$. The map $O[[Y]]\to A$ is
finite free and faithfully flat, with basis $1,v,\ldots,v^{p-1}$:
the preparation theorem applied to $P_1(v)-Y$ in $O[[Y]][[v]]$
gives a monic distinguished polynomial of degree $p$.
Thus membership in a scalar-uniformizer power modulo $P_2^y$
for series in $Y$ can be tested in $O[[Y]]$. Indeed $P_2$ is
itself a series in $Y$, of the form $p\varepsilon(Y)+U(Y)Y^p$
with both coefficients units. This follows by applying Frobenius
to $P_1=p+(P_0-p)P_0^{p-1}U_0(P_0)$ in the preceding letter.
Put $Z=U(Y)Y^p/\varepsilon(Y)$, so $P_2=\varepsilon(Y)(p+Z)$.
Again $O[[Y]]$ is finite free over $O[[Z]]$ with a basis
containing $1$. The content of $Y^{b_1}$ is bounded above by
that of $Y^{pd}$, which is a unit times $Z^d$. In
$O[[Z]]/((p+Z)^y)$ use the basis $(p+Z)^j$, $0\le j<y$.
The coefficient at $j=y-1$ is
$\binom d{y-1}(-p)^{d-y+1}$. Its valuation is exactly $K$.
These tests work in a ramified coefficient field with $v_p(p)=1$:
one may equivalently test all integer powers of its uniformizer.
It follows that $p^Kx_2$ is integral. The indicated integral
numerator for $f_2$ belongs to $P_2^yA$ after inverting $p$,
hence already over $O$ by distinguished division. Subtracting
the same constant in the $P_0$ congruence and clearing the
remaining scalar denominator gives the last assertion.

In the nonzero-cost case $py\le b_1<p(y+\eta)$ gives
$d\le y+\eta$, while $v_p\binom d{y-1}\le\log_p b_1$.
This is a band place: if $b_*$ is the thick threshold, then
$b_*\le b_1\le b_*+\eta$, so
$|b_*-p s_{c-2}|\le p(S+\eta)+\eta$.
In fact its negative part is at most $\eta$ and its positive
part at most $p(S+\eta)$, each allowed by the native band
quantity with constant $C_b\ge2$. The adjacent thick letter
is $c-1$. Since $b_1\le2h$, the cost is at most
$(\eta+2+\log_p(2h))\,\mathbf1_{\rm band}$. This proves
the native three-letter scalar descent with its band guard.
% END K83J1

% BEGIN K83J2
\paragraph{K83J2. The lower-level constant after descent.}
For the same $f_2,\kappa,X_0$, every nonzero $X_0$ satisfies
\[
 v_p(X_0)\ge \kappa+y-b_1+B/p-B/p^2-2-\log_p(2h).
\]
Write $Q=P_0=v+p$. In $A/(Q^B)$ its defining equation gives
$X_0 P_1^{b_1}\equiv p^\kappa P_2^y a$ for some $a\in A$.
As a series in $Q$, the reduction of $P_2$ modulo $p$ is a
unit times $Q^{p^2}$, by iterating $\varphi(v)\sim vQ^{p-1}$.
Thus $P_2=pE+Q^{p^2}V$ for two integral series. The surviving
terms in the $y$th power modulo $Q^B$ use fewer than $B/p^2$
copies of the second summand; they all have valuation at least
$y-B/p^2$. The right side therefore has content at least
$\kappa+y-B/p^2$.

The content of $P_1^{b_1}$ in this same finite free quotient
is at most $b_1-B/p+1+\log_p b_1$. Indeed $P_1=p\varepsilon+
Q^p U$ with units $\varepsilon,U$; apply T83A2's
first-offending-coefficient argument at $j=\lfloor(B-1)/p\rfloor$.
Here $1\le B\le p b_1$ because both terminal letters are thick,
$S,\eta$ are each at most $b_*/20$ on the late mixed family, and
$p\ge11$. Subtracting these two content bounds proves the displayed
inequality, with an extra unit of slack, for the exact native
normalization. If $X_0=0$ its valuation is infinite.
% END K83J2

% BEGIN K83J3
\paragraph{K83J3. A lift of the separated constant on the original map.}
Keep the target horizontal basis of U83M3 and any chosen
coordinate $k$. For any $x\in O$, the native scalar $x$ in that
coordinate has a source $\mathrm{LiftsAt}$ witness of exponent
$R+\sigma$, where $R$ and $\sigma$ are the fixed U83M1 witnesses.
In particular this lifts $X_2$ from K83J1. At the preceding cut
$c-3$ lift $p^Rx$ times the target basis vector through $F$ in
the invariant Wach lattice by U83M1. Its specialization is an
integral source residue vector $a$. Use a source horizontal basis
and U83M6's forward matrix to realize its constant coordinates
at $c$ at cost $p^\sigma$. Call this vector $w_1$.
Its generator difference lies in $vN_c(T_1)$ since the scalar
coordinates are constant. The vector is invariant and has finite
cyclotomic poles by the two translated height bounds. Its image
belongs to the required localized target Frobenius span since the
map preserves the integral translated span.

To check the native representing congruence, in source and target
bases let $M_F(v)$ be the integral matrix of $F$ at $c-3$ and
$U_1,U_2$ the normalized horizontal matrices. The formal equation
$M_F(v)U_1(v)=U_2(v)M_F(0)$ follows by uniqueness of the
invariant series with its given constant. Apply $\varphi^3$
to its truncation at $v^M$. The omitted terms vanish modulo
$v\Lambda^2$ by U83M6. Hence the target representing coordinates
of $Fw_1$ are exactly $p^{R+\sigma}x$ in coordinate $k$ and
zero elsewhere. The additional vector in $vN_c(T_2)$ in the
definition of $\mathrm{LiftsAt}$ may be taken zero.
This realizes the native constant-term lift before any further
multiplication. For this sign and the guard $b_1<p(y+\eta)$,
the still-required nonconstant operation is to lift $t_3f_2$
on the same source with exponent $e\le6U$ and scalar
$v_p(t_3)\le14U$, with $U$ as in U83M1. K83J2 supplies its
lower-level content estimate.
% END K83J3
